\section{Conclusion}
\label{sec:conclusion}

In this work, we introduced \textbf{Sparsity-Guided Task Arithmetic (SG-TA)}, a simple, deterministic, and robust weight-space regularization framework for model merging. Rather than relying on complex stochastic masking or multi-stage sign election protocols, SG-TA decouples the sparsification process by applying absolute magnitude-based binary masks to task vectors before merging. 

Through exhaustive, multi-axial empirical sweeps across 4 visual classification benchmarks (MNIST, FashionMNIST, CIFAR-10, SVHN) using a Vision Transformer backbone, we validated our core hypotheses. First, magnitude-based binary masking acts as an exceptional spatial regularizer, mitigating the catastrophic representation collapse inherent to naive Task Arithmetic. Second, Global Quantile (GQ) masking significantly outperforms Layer-wise Quantile (LQ) masking (achieving a joint mean accuracy of $61.40\%$ vs. $57.81\%$), demonstrating that global budget flexibility is vital for preserving active updates in task-sensitive layers. Third, simple deterministic magnitude masks outperform state-of-the-art baselines like TIES-Merging and DARE-Merging on Vision Transformers, highlighting that structural determinism is superior to stochastic dropout and complex majority sign heuristics.

Adhering to our empirical values, we showed that the combination of simple design and rigorous calibration (via Offline Few-Shot Validation Tuning) provides a robust, lightweight, and effective path forward for model merging. Future work will investigate scaling SG-TA to large language models (LLMs), exploring unsupervised calibration alternatives to bypass labeled sample requirements without triggering test-time adaptation collapse, and integrating first-order Fisher-weighted gates or soft/elastic regularizers to bridge the absolute performance gap.
